% tp2.tex — Trabalho Prático II
\documentclass[12pt]{article}

\usepackage[utf8]{inputenc}
\usepackage[T1]{fontenc}
\usepackage{lmodern}
\usepackage[portuguese]{babel}
\usepackage{geometry}
\geometry{a4paper, margin=2.5cm}

\usepackage{amsmath, amssymb, amsfonts}
\usepackage{enumitem}
\usepackage{xcolor}
\usepackage{hyperref}
\usepackage{listings}
\lstset{
  basicstyle=\ttfamily\small,
  frame=single,
  breaklines=true,
  tabsize=2
}

\title{Trabalho Prático II — Pesquisa Aplicada\\Ataques por Inversão de Gradiente}
\author{Disciplina de Otimização Não Linear}
\date{\today}

\begin{document}
\maketitle

\section*{Objetivo}

Estender o conhecimento adquirido no TP1 propondo, desenvolvendo e avaliando uma \textbf{pesquisa aplicada} envolvendo \emph{ataques por inversão de gradiente} (Gradient Inversion Attacks, GIA). O trabalho pode ter natureza empírica, teórica ou metodológica.

\section*{Possibilidades de Pesquisa}

Você pode escolher uma das seguintes linhas de investigação (ou propor outra, mediante aprovação do professor):

\begin{enumerate}[label=\textbf{\arabic*.}, nosep]
  \item \textbf{Aplicação em tarefas diferentes:}
        \begin{itemize}
          \item Aplicar GIA em outros conjuntos de dados (ex: CIFAR-10, Fashion-MNIST).
          \item Explorar GIA em tarefas não classificatórias (ex: regressão).
        \end{itemize}
  \item \textbf{Variação arquitetural:}
        \begin{itemize}
          \item Investigar o impacto da arquitetura do modelo (ex: redes convolucionais, redes mais profundas).
          \item GIA em modelos com regularização (Dropout, BatchNorm etc.).
        \end{itemize}
  \item \textbf{Robustez e defesa:}
        \begin{itemize}
          \item Avaliar a eficácia de defesas conhecidas (ex: clipping, ruído de privacidade diferencial).
          \item Propor novas técnicas de defesa contra GIA.
        \end{itemize}
  \item \textbf{Aprimoramento do ataque:}
        \begin{itemize}
          \item Substituir a função de perda ou adicionar termos auxiliares (TV-loss, latentes de GANs).
          \item Propor inicializações ou otimizações alternativas.
        \end{itemize}
  \item \textbf{Estudo comparativo:}
        \begin{itemize}
          \item Comparar empiricamente diferentes variantes (DLG, iDLG, DLG-FB).
          \item Analisar a influência de diferentes otimizadores (Adam, L-BFGS, SGD).
        \end{itemize}
\end{enumerate}

\section*{Entregáveis (100 pts)}

\begin{itemize}[nosep]
  \item \textbf{Artigo (100 pts)} — entre 10 e 15 páginas, escrito em português ou inglês, usando o template disponível em:
    \begin{center}
      \url{https://github.com/rcpsilva/BCC405-PCC179_NonLinearOptimization/tree/main/2025-1}
    \end{center}
\end{itemize}

O artigo deve conter:

\begin{enumerate}[label=\textbf{\arabic*.}, nosep]
  \item \textbf{Introdução} — contextualização do problema, motivação e objetivos do trabalho.
  \item \textbf{Trabalhos Relacionados} — breve revisão dos principais trabalhos da literatura.
  \item \textbf{Metodologia} — descrição da hipótese ou proposta investigada; detalhes do modelo, dados, métricas e configuração experimental.
  \item \textbf{Resultados e Discussão} — análise quantitativa e qualitativa; comparação de abordagens e visualizações dos resultados.
  \item \textbf{Conclusão} — principais achados, limitações e sugestões para trabalhos futuros.
\end{enumerate}

\section*{Critérios de Avaliação}

\begin{itemize}[nosep]
  \item \textbf{15 pts} — Clareza na definição da proposta e motivação.
  \item \textbf{15 pts} — Fundamentação teórica e revisão de literatura.
  \item \textbf{20 pts} — Qualidade metodológica e reprodutibilidade.
  \item \textbf{20 pts} — Profundidade da análise de resultados.
  \item \textbf{15 pts} — Escrita e organização do artigo.
  \item \textbf{15 pts} — Criatividade/inovação e complexidade da proposta.
\end{itemize}

\section*{Datas Importantes}

\begin{itemize}[nosep]
  \item \textbf{Entrega do artigo:} 24/08/2025 - 23h59.
\end{itemize}

\section*{Observações Finais}

Dúvidas ou propostas fora do escopo podem ser discutidas em horário de atendimento ou via e-mail. Encorajamos a criatividade, o pensamento crítico e a investigação rigorosa.

\end{document}
